%%%%%%%%%%%%%%%%%%%%%%%%%%%%%%%%%%%%%%%%%%%%%%%%%%%%%%%%%%%%
% 7.11 MEASURE-THEORETIC PROBABILITY
% The Composition of Advanced Mathematics
%%%%%%%%%%%%%%%%%%%%%%%%%%%%%%%%%%%%%%%%%%%%%%%%%%%%%%%%%%%%

\section{Measure-Theoretic Probability}
\label{sec:measure-theoretic-probability}

\prerequisite{
Elementary set theory, functions, real analysis, sequences of functions,
measure theory, Lebesgue integration, random variables, expectation,
joint distributions, laws of large numbers, and central limit
theorems.
}

\learningobjectives{
By the end of this section, the reader should be able to:
\begin{itemize}
    \item define a probability space;
    \item explain the roles of the sample space, sigma-algebra, and
          probability measure;
    \item define measurable events;
    \item interpret Kolmogorov's axioms measure-theoretically;
    \item derive standard probability identities from measure axioms;
    \item define random variables as measurable functions;
    \item define distributions as pushforward probability measures;
    \item distinguish discrete, absolutely continuous, and singular laws;
    \item define expectation as a Lebesgue integral;
    \item work with indicator functions and simple random variables;
    \item apply monotone and dominated convergence theorems to
          expectations;
    \item use Fatou's lemma in probabilistic settings;
    \item understand convergence almost surely, in probability, in
          distribution, and in \(L^p\);
    \item understand relationships among convergence modes;
    \item define conditional expectation with respect to a sigma-algebra;
    \item use the tower property;
    \item understand conditional probability as conditional expectation
          of an indicator;
    \item define independence of sigma-algebras and random variables;
    \item understand product probability spaces;
    \item define \(L^p\) spaces of random variables;
    \item interpret variance geometrically in \(L^2\);
    \item recognize conditional expectation as an \(L^2\) projection;
    \item state the Radon--Nikodym theorem in probabilistic language;
    \item understand regular conditional probability conceptually;
    \item use Borel--Cantelli lemmas rigorously;
    \item connect measure-theoretic probability with stochastic processes,
          martingales, stochastic calculus, and modern statistics.
\end{itemize}
}

%%%%%%%%%%%%%%%%%%%%%%%%%%%%%%%%%%%%%%%%%%%%%%%%%%%%%%%%%%%%
\subsection{Why Probability Needs Measure Theory}
\label{subsec:why-probability-measure-theory}
%%%%%%%%%%%%%%%%%%%%%%%%%%%%%%%%%%%%%%%%%%%%%%%%%%%%%%%%%%%%

Elementary probability often begins with finite or countable sample
spaces.

For example,

\[
\Omega
=
\{1,2,3,4,5,6\}
\]

for a die roll.

In continuous probability, however, the sample space may be

\[
\Omega
=
\mathbb{R},
\]

or an infinite-dimensional space of functions or sequences.

Probability must then answer questions such as:

\[
P(X\in A)
\]

for complicated sets \(A\).

Measure theory provides the rigorous framework needed to define such
probabilities consistently.

%%%%%%%%%%%%%%%%%%%%%%%%%%%%%%%%%%%%%%%%%%%%%%%%%%%%%%%%%%%%
\subsection{Probability Space}
\label{subsec:probability-space}
%%%%%%%%%%%%%%%%%%%%%%%%%%%%%%%%%%%%%%%%%%%%%%%%%%%%%%%%%%%%

\begin{definitionbox}{Probability Space}
A probability space is a triple

\[
\boxed{
(\Omega,\mathcal{F},P),
}
\]

where:

\begin{itemize}
    \item \(\Omega\) is the sample space;
    \item \(\mathcal{F}\) is a sigma-algebra of subsets of \(\Omega\);
    \item \(P\) is a probability measure on \((\Omega,\mathcal{F})\).
\end{itemize}
\end{definitionbox}

The symbol

\[
\Omega
\]

is uppercase Greek omega.

The symbol

\[
\mathcal{F}
\]

is commonly used for the collection of measurable events.

%%%%%%%%%%%%%%%%%%%%%%%%%%%%%%%%%%%%%%%%%%%%%%%%%%%%%%%%%%%%
\subsection{Sigma-Algebras}
\label{subsec:sigma-algebras-probability}
%%%%%%%%%%%%%%%%%%%%%%%%%%%%%%%%%%%%%%%%%%%%%%%%%%%%%%%%%%%%

\begin{definitionbox}{Sigma-Algebra}
A collection

\[
\mathcal{F}
\subseteq
2^\Omega
\]

is a sigma-algebra if:

\[
\boxed{
\Omega\in\mathcal{F},
}
\]

if

\[
A\in\mathcal{F},
\]

then

\[
\boxed{
A^c\in\mathcal{F},
}
\]

and if

\[
A_1,A_2,\ldots\in\mathcal{F},
\]

then

\[
\boxed{
\bigcup_{n=1}^{\infty}A_n
\in
\mathcal{F}.
}
\]
\end{definitionbox}

The prefix ``sigma'' indicates closure under countable operations.

%%%%%%%%%%%%%%%%%%%%%%%%%%%%%%%%%%%%%%%%%%%%%%%%%%%%%%%%%%%%
\subsection{Measurable Events}
\label{subsec:measurable-events}
%%%%%%%%%%%%%%%%%%%%%%%%%%%%%%%%%%%%%%%%%%%%%%%%%%%%%%%%%%%%

An event is measurable if it belongs to the sigma-algebra.

Thus,

\[
\boxed{
A\in\mathcal{F}
}
\]

means \(A\) is an event to which the probability measure assigns a
probability.

Not every subset of an uncountable sample space must be measurable.

This distinction becomes important in advanced analysis.

%%%%%%%%%%%%%%%%%%%%%%%%%%%%%%%%%%%%%%%%%%%%%%%%%%%%%%%%%%%%
\subsection{Basic Sigma-Algebra Consequences}
\label{subsec:sigma-algebra-consequences}
%%%%%%%%%%%%%%%%%%%%%%%%%%%%%%%%%%%%%%%%%%%%%%%%%%%%%%%%%%%%

From the axioms, one obtains

\[
\varnothing\in\mathcal{F}.
\]

Indeed,

\[
\varnothing
=
\Omega^c.
\]

Sigma-algebras are also closed under countable intersections because

\[
\bigcap_{n=1}^{\infty}A_n
=
\left(
\bigcup_{n=1}^{\infty}A_n^c
\right)^c.
\]

Thus,

\[
\boxed{
A_n\in\mathcal{F}
\Longrightarrow
\bigcap_{n=1}^{\infty}A_n\in\mathcal{F}.
}
\]

%%%%%%%%%%%%%%%%%%%%%%%%%%%%%%%%%%%%%%%%%%%%%%%%%%%%%%%%%%%%
\subsection{Generated Sigma-Algebras}
\label{subsec:generated-sigma-algebras}
%%%%%%%%%%%%%%%%%%%%%%%%%%%%%%%%%%%%%%%%%%%%%%%%%%%%%%%%%%%%

Given a collection

\[
\mathcal{C}
\]

of subsets of \(\Omega\), the sigma-algebra generated by \(\mathcal{C}\)
is written

\[
\boxed{
\sigma(\mathcal{C}).
}
\]

It is the smallest sigma-algebra containing every set in
\(\mathcal{C}\).

Generated sigma-algebras allow complicated measurable structures to be
built from simple sets.

%%%%%%%%%%%%%%%%%%%%%%%%%%%%%%%%%%%%%%%%%%%%%%%%%%%%%%%%%%%%
\subsection{Borel Sigma-Algebra}
\label{subsec:borel-sigma-algebra-probability}
%%%%%%%%%%%%%%%%%%%%%%%%%%%%%%%%%%%%%%%%%%%%%%%%%%%%%%%%%%%%

The Borel sigma-algebra on \(\mathbb{R}\) is written

\[
\boxed{
\mathcal{B}(\mathbb{R}).
}
\]

It is generated by the open subsets of \(\mathbb{R}\).

Equivalently, it may be generated by intervals of forms such as

\[
(-\infty,x],
\]

\[
(a,b),
\]

or

\[
[a,b].
\]

Most ordinary real-valued probability distributions are defined on
Borel sets.

%%%%%%%%%%%%%%%%%%%%%%%%%%%%%%%%%%%%%%%%%%%%%%%%%%%%%%%%%%%%
\subsection{Probability Measure}
\label{subsec:probability-measure}
%%%%%%%%%%%%%%%%%%%%%%%%%%%%%%%%%%%%%%%%%%%%%%%%%%%%%%%%%%%%

\begin{definitionbox}{Probability Measure}
A probability measure is a function

\[
P:\mathcal{F}\to[0,1]
\]

satisfying:

\[
\boxed{
P(\Omega)=1,
}
\]

and for pairwise disjoint events

\[
A_1,A_2,\ldots,
\]

\[
\boxed{
P
\left(
\bigcup_{n=1}^{\infty}A_n
\right)
=
\sum_{n=1}^{\infty}P(A_n).
}
\]
\end{definitionbox}

The second condition is called countable additivity.

%%%%%%%%%%%%%%%%%%%%%%%%%%%%%%%%%%%%%%%%%%%%%%%%%%%%%%%%%%%%
\subsection{Kolmogorov Axioms}
\label{subsec:kolmogorov-axioms-measure}
%%%%%%%%%%%%%%%%%%%%%%%%%%%%%%%%%%%%%%%%%%%%%%%%%%%%%%%%%%%%

The usual Kolmogorov axioms can be written as:

\[
\boxed{
P(A)\geq0,
}
\]

\[
\boxed{
P(\Omega)=1,
}
\]

and for pairwise disjoint events,

\[
\boxed{
P
\left(
\bigcup_{n=1}^{\infty}A_n
\right)
=
\sum_{n=1}^{\infty}P(A_n).
}
\]

These are precisely the axioms of a finite measure normalized so that
the entire space has measure \(1\).

%%%%%%%%%%%%%%%%%%%%%%%%%%%%%%%%%%%%%%%%%%%%%%%%%%%%%%%%%%%%
\subsection{Probability as a Measure}
\label{subsec:probability-as-measure}
%%%%%%%%%%%%%%%%%%%%%%%%%%%%%%%%%%%%%%%%%%%%%%%%%%%%%%%%%%%%

Probability theory is therefore a special case of measure theory.

The correspondence is:

\[
\boxed{
\text{measurable set}
\longleftrightarrow
\text{event},
}
\]

\[
\boxed{
\text{measure}
\longleftrightarrow
\text{probability},
}
\]

and

\[
\boxed{
\text{measurable function}
\longleftrightarrow
\text{random variable}.
}
\]

This is the central structural viewpoint of modern probability.

%%%%%%%%%%%%%%%%%%%%%%%%%%%%%%%%%%%%%%%%%%%%%%%%%%%%%%%%%%%%
\subsection{Consequences of the Probability Axioms}
\label{subsec:probability-measure-consequences}
%%%%%%%%%%%%%%%%%%%%%%%%%%%%%%%%%%%%%%%%%%%%%%%%%%%%%%%%%%%%

Since

\[
\Omega
=
\Omega\cup\varnothing
\]

with disjoint sets,

\[
P(\Omega)
=
P(\Omega)+P(\varnothing).
\]

Therefore,

\[
\boxed{
P(\varnothing)=0.
}
\]

If

\[
A\subseteq B,
\]

then

\[
B
=
A
\cup
(B\setminus A)
\]

disjointly.

Thus,

\[
\boxed{
P(A)\leq P(B).
}
\]

This is monotonicity.

%%%%%%%%%%%%%%%%%%%%%%%%%%%%%%%%%%%%%%%%%%%%%%%%%%%%%%%%%%%%
\subsection{Complement Formula}
\label{subsec:measure-complement-formula}
%%%%%%%%%%%%%%%%%%%%%%%%%%%%%%%%%%%%%%%%%%%%%%%%%%%%%%%%%%%%

Since

\[
A\cup A^c=\Omega
\]

and the union is disjoint,

\[
P(A)+P(A^c)=1.
\]

Therefore,

\[
\boxed{
P(A^c)=1-P(A).
}
\]

%%%%%%%%%%%%%%%%%%%%%%%%%%%%%%%%%%%%%%%%%%%%%%%%%%%%%%%%%%%%
\subsection{Continuity from Below}
\label{subsec:continuity-from-below-probability}
%%%%%%%%%%%%%%%%%%%%%%%%%%%%%%%%%%%%%%%%%%%%%%%%%%%%%%%%%%%%

Suppose

\[
A_1\subseteq A_2\subseteq\cdots.
\]

Then

\[
\boxed{
P
\left(
\bigcup_{n=1}^{\infty}A_n
\right)
=
\lim_{n\to\infty}P(A_n).
}
\]

This is called continuity from below.

It follows from countable additivity.

%%%%%%%%%%%%%%%%%%%%%%%%%%%%%%%%%%%%%%%%%%%%%%%%%%%%%%%%%%%%
\subsection{Continuity from Above}
\label{subsec:continuity-from-above-probability}
%%%%%%%%%%%%%%%%%%%%%%%%%%%%%%%%%%%%%%%%%%%%%%%%%%%%%%%%%%%%

Suppose

\[
A_1\supseteq A_2\supseteq\cdots.
\]

Then because \(P(A_1)<\infty\),

\[
\boxed{
P
\left(
\bigcap_{n=1}^{\infty}A_n
\right)
=
\lim_{n\to\infty}P(A_n).
}
\]

This is continuity from above.

%%%%%%%%%%%%%%%%%%%%%%%%%%%%%%%%%%%%%%%%%%%%%%%%%%%%%%%%%%%%
\subsection{Worked Example: Increasing Events}
\label{subsec:worked-increasing-events}
%%%%%%%%%%%%%%%%%%%%%%%%%%%%%%%%%%%%%%%%%%%%%%%%%%%%%%%%%%%%

\begin{workedexamplebox}{An Increasing Sequence of Events}

Let

\[
A_n
=
\left[
-\!n,n
\right]
\]

on a real-valued probability space.

Then

\[
A_1\subseteq A_2\subseteq\cdots
\]

and

\[
\bigcup_{n=1}^{\infty}A_n
=
\mathbb{R}.
\]

Therefore,

\[
P(A_n)\to P(\mathbb{R}).
\]

\begin{answerbox}
\[
\boxed{
P([-n,n])\to1.
}
\]
\end{answerbox}

\end{workedexamplebox}

%%%%%%%%%%%%%%%%%%%%%%%%%%%%%%%%%%%%%%%%%%%%%%%%%%%%%%%%%%%%
\subsection{Random Variables as Measurable Functions}
\label{subsec:random-variables-measurable-functions}
%%%%%%%%%%%%%%%%%%%%%%%%%%%%%%%%%%%%%%%%%%%%%%%%%%%%%%%%%%%%

\begin{definitionbox}{Random Variable}
A real-valued random variable on

\[
(\Omega,\mathcal{F},P)
\]

is a measurable function

\[
\boxed{
X:
(\Omega,\mathcal{F})
\to
(\mathbb{R},\mathcal{B}(\mathbb{R})).
}
\]
\end{definitionbox}

Measurability means

\[
\boxed{
X^{-1}(B)
\in
\mathcal{F}
}
\]

for every Borel set

\[
B\in\mathcal{B}(\mathbb{R}).
\]

%%%%%%%%%%%%%%%%%%%%%%%%%%%%%%%%%%%%%%%%%%%%%%%%%%%%%%%%%%%%
\subsection{Why Measurability Matters}
\label{subsec:why-rv-measurability}
%%%%%%%%%%%%%%%%%%%%%%%%%%%%%%%%%%%%%%%%%%%%%%%%%%%%%%%%%%%%

To write

\[
P(X\in B),
\]

we really mean

\[
P
\left(
\{\omega:X(\omega)\in B\}
\right).
\]

The set inside the probability is

\[
X^{-1}(B).
\]

Measurability ensures that this preimage is an event in
\(\mathcal{F}\).

Thus probability statements involving \(X\) are well-defined.

%%%%%%%%%%%%%%%%%%%%%%%%%%%%%%%%%%%%%%%%%%%%%%%%%%%%%%%%%%%%
\subsection{Equivalent Measurability Test}
\label{subsec:rv-measurability-test}
%%%%%%%%%%%%%%%%%%%%%%%%%%%%%%%%%%%%%%%%%%%%%%%%%%%%%%%%%%%%

For a real-valued function \(X\), it is enough to verify that

\[
\boxed{
\{X\leq x\}
\in
\mathcal{F}
}
\]

for every real \(x\).

Equivalently, one may use events such as

\[
\{X<x\},
\]

\[
\{X>x\},
\]

or

\[
\{X\geq x\}.
\]

These sets generate the Borel sigma-algebra.

%%%%%%%%%%%%%%%%%%%%%%%%%%%%%%%%%%%%%%%%%%%%%%%%%%%%%%%%%%%%
\subsection{Distribution as a Pushforward Measure}
\label{subsec:distribution-pushforward}
%%%%%%%%%%%%%%%%%%%%%%%%%%%%%%%%%%%%%%%%%%%%%%%%%%%%%%%%%%%%

A random variable \(X\) induces a probability measure on
\(\mathbb{R}\).

\begin{definitionbox}{Distribution of a Random Variable}
The distribution, or law, of \(X\) is the measure

\[
\boxed{
P_X(B)
=
P(X\in B)
=
P(X^{-1}(B)).
}
\]
\end{definitionbox}

This is the pushforward of \(P\) through \(X\).

One writes

\[
\boxed{
P_X
=
P\circ X^{-1}.
}
\]

%%%%%%%%%%%%%%%%%%%%%%%%%%%%%%%%%%%%%%%%%%%%%%%%%%%%%%%%%%%%
\subsection{Same Distribution, Different Probability Spaces}
\label{subsec:same-law-different-spaces}
%%%%%%%%%%%%%%%%%%%%%%%%%%%%%%%%%%%%%%%%%%%%%%%%%%%%%%%%%%%%

Two random variables may be defined on completely different probability
spaces and still have the same distribution.

We write

\[
\boxed{
X\overset{d}{=}Y
}
\]

if

\[
P_X=P_Y.
\]

Thus a distribution describes the probabilistic behavior of the values,
not the underlying sample-space construction.

%%%%%%%%%%%%%%%%%%%%%%%%%%%%%%%%%%%%%%%%%%%%%%%%%%%%%%%%%%%%
\subsection{Distribution Function Measure-Theoretically}
\label{subsec:cdf-measure-theoretic}
%%%%%%%%%%%%%%%%%%%%%%%%%%%%%%%%%%%%%%%%%%%%%%%%%%%%%%%%%%%%

For a real-valued random variable,

\[
F_X(x)
=
P(X\leq x).
\]

Using the pushforward measure,

\[
\boxed{
F_X(x)
=
P_X((-\infty,x]).
}
\]

Thus the CDF is simply the distribution measure evaluated on a
particular family of Borel sets.

%%%%%%%%%%%%%%%%%%%%%%%%%%%%%%%%%%%%%%%%%%%%%%%%%%%%%%%%%%%%
\subsection{Discrete Laws}
\label{subsec:discrete-laws-measure}
%%%%%%%%%%%%%%%%%%%%%%%%%%%%%%%%%%%%%%%%%%%%%%%%%%%%%%%%%%%%

If \(X\) takes countably many values

\[
x_1,x_2,\ldots,
\]

then

\[
P_X
\]

is concentrated on those points.

For any Borel set \(B\),

\[
\boxed{
P_X(B)
=
\sum_{x_i\in B}
P(X=x_i).
}
\]

Such a law is called discrete.

%%%%%%%%%%%%%%%%%%%%%%%%%%%%%%%%%%%%%%%%%%%%%%%%%%%%%%%%%%%%
\subsection{Absolutely Continuous Laws}
\label{subsec:absolutely-continuous-laws}
%%%%%%%%%%%%%%%%%%%%%%%%%%%%%%%%%%%%%%%%%%%%%%%%%%%%%%%%%%%%

A probability law \(P_X\) is absolutely continuous with respect to
Lebesgue measure if there exists a nonnegative function \(f_X\) such
that

\[
\boxed{
P_X(B)
=
\int_B
f_X(x)\,dx
}
\]

for every Borel set \(B\).

The function \(f_X\) is the probability density.

%%%%%%%%%%%%%%%%%%%%%%%%%%%%%%%%%%%%%%%%%%%%%%%%%%%%%%%%%%%%
\subsection{Density as a Radon--Nikodym Derivative}
\label{subsec:density-radon-nikodym}
%%%%%%%%%%%%%%%%%%%%%%%%%%%%%%%%%%%%%%%%%%%%%%%%%%%%%%%%%%%%

Measure-theoretically,

\[
\boxed{
f_X
=
\frac{dP_X}{dx},
}
\]

where \(dx\) denotes Lebesgue measure.

More precisely,

\[
\boxed{
f_X
=
\frac{dP_X}{d\lambda},
}
\]

where \(\lambda\) is Lebesgue measure.

Thus a PDF is a Radon--Nikodym derivative.

%%%%%%%%%%%%%%%%%%%%%%%%%%%%%%%%%%%%%%%%%%%%%%%%%%%%%%%%%%%%
\subsection{Singular Distributions}
\label{subsec:singular-distributions}
%%%%%%%%%%%%%%%%%%%%%%%%%%%%%%%%%%%%%%%%%%%%%%%%%%%%%%%%%%%%

Not every continuous CDF has a density.

A probability measure may be continuous but singular with respect to
Lebesgue measure.

The classic example is the Cantor distribution.

It assigns probability \(1\) to the Cantor set, which has Lebesgue
measure \(0\).

Thus probability distributions can be:

\[
\boxed{
\text{discrete},
}
\]

\[
\boxed{
\text{absolutely continuous},
}
\]

or

\[
\boxed{
\text{singular}.
}
\]

Mixtures of these types are also possible.

%%%%%%%%%%%%%%%%%%%%%%%%%%%%%%%%%%%%%%%%%%%%%%%%%%%%%%%%%%%%
\subsection{Indicator Functions}
\label{subsec:indicator-functions-measure}
%%%%%%%%%%%%%%%%%%%%%%%%%%%%%%%%%%%%%%%%%%%%%%%%%%%%%%%%%%%%

For an event \(A\),

\[
\mathbf{1}_A(\omega)
=
\begin{cases}
1,&\omega\in A,\\
0,&\omega\notin A.
\end{cases}
\]

Then

\[
\boxed{
\mathbb{E}[\mathbf{1}_A]
=
P(A).
}
\]

This identity links integration directly to probability.

%%%%%%%%%%%%%%%%%%%%%%%%%%%%%%%%%%%%%%%%%%%%%%%%%%%%%%%%%%%%
\subsection{Simple Random Variables}
\label{subsec:simple-random-variables}
%%%%%%%%%%%%%%%%%%%%%%%%%%%%%%%%%%%%%%%%%%%%%%%%%%%%%%%%%%%%

A simple random variable takes only finitely many values.

It can be written

\[
\boxed{
X
=
\sum_{k=1}^{m}
a_k
\mathbf{1}_{A_k},
}
\]

where the \(A_k\) are measurable events.

For nonnegative simple \(X\),

\[
\boxed{
\mathbb{E}[X]
=
\sum_{k=1}^{m}
a_kP(A_k).
}
\]

This is the starting point for defining the Lebesgue integral.

%%%%%%%%%%%%%%%%%%%%%%%%%%%%%%%%%%%%%%%%%%%%%%%%%%%%%%%%%%%%
\subsection{Expectation as a Lebesgue Integral}
\label{subsec:expectation-lebesgue-integral}
%%%%%%%%%%%%%%%%%%%%%%%%%%%%%%%%%%%%%%%%%%%%%%%%%%%%%%%%%%%%

For a measurable random variable \(X\),

\[
\boxed{
\mathbb{E}[X]
=
\int_\Omega
X(\omega)\,dP(\omega),
}
\]

whenever the integral is defined.

This is the fundamental measure-theoretic definition of expectation.

It includes discrete sums and continuous density integrals as special
cases.

%%%%%%%%%%%%%%%%%%%%%%%%%%%%%%%%%%%%%%%%%%%%%%%%%%%%%%%%%%%%
\subsection{Expectation Using the Distribution Measure}
\label{subsec:expectation-law-integral}
%%%%%%%%%%%%%%%%%%%%%%%%%%%%%%%%%%%%%%%%%%%%%%%%%%%%%%%%%%%%

The law of the unconscious statistician becomes

\[
\boxed{
\mathbb{E}[g(X)]
=
\int_{\mathbb{R}}
g(x)\,dP_X(x),
}
\]

provided the integral is defined.

If \(P_X\) has density \(f_X\),

\[
\boxed{
\mathbb{E}[g(X)]
=
\int_{\mathbb{R}}
g(x)f_X(x)\,dx.
}
\]

If \(X\) is discrete,

\[
\boxed{
\mathbb{E}[g(X)]
=
\sum_x
g(x)P(X=x).
}
\]

%%%%%%%%%%%%%%%%%%%%%%%%%%%%%%%%%%%%%%%%%%%%%%%%%%%%%%%%%%%%
\subsection{Positive and Negative Parts}
\label{subsec:positive-negative-parts}
%%%%%%%%%%%%%%%%%%%%%%%%%%%%%%%%%%%%%%%%%%%%%%%%%%%%%%%%%%%%

Define

\[
\boxed{
X^+
=
\max(X,0)
}
\]

and

\[
\boxed{
X^-
=
\max(-X,0).
}
\]

Then

\[
X
=
X^+-X^-,
\]

and

\[
|X|
=
X^++X^-.
\]

Expectation is defined through

\[
\boxed{
\mathbb{E}[X]
=
\mathbb{E}[X^+]
-
\mathbb{E}[X^-]
}
\]

provided the expression is not of the indeterminate form

\[
\infty-\infty.
\]

%%%%%%%%%%%%%%%%%%%%%%%%%%%%%%%%%%%%%%%%%%%%%%%%%%%%%%%%%%%%
\subsection{Integrable Random Variables}
\label{subsec:integrable-random-variables}
%%%%%%%%%%%%%%%%%%%%%%%%%%%%%%%%%%%%%%%%%%%%%%%%%%%%%%%%%%%%

A random variable is integrable if

\[
\boxed{
\mathbb{E}[|X|]<\infty.
}
\]

This implies that both

\[
\mathbb{E}[X^+]
\]

and

\[
\mathbb{E}[X^-]
\]

are finite.

Therefore \(\mathbb{E}[X]\) exists as a finite real number.

%%%%%%%%%%%%%%%%%%%%%%%%%%%%%%%%%%%%%%%%%%%%%%%%%%%%%%%%%%%%
\subsection{Almost Everywhere and Almost Surely}
\label{subsec:ae-as}
%%%%%%%%%%%%%%%%%%%%%%%%%%%%%%%%%%%%%%%%%%%%%%%%%%%%%%%%%%%%

In measure theory, a property holds almost everywhere if it fails only
on a set of measure zero.

In probability, the corresponding phrase is almost surely.

Thus,

\[
\boxed{
\text{almost everywhere with respect to }P
=
\text{almost surely}.
}
\]

For example,

\[
X=Y
\quad\text{a.s.}
\]

means

\[
\boxed{
P(X=Y)=1.
}
\]

%%%%%%%%%%%%%%%%%%%%%%%%%%%%%%%%%%%%%%%%%%%%%%%%%%%%%%%%%%%%
\subsection{Monotone Convergence Theorem}
\label{subsec:mct-probability}
%%%%%%%%%%%%%%%%%%%%%%%%%%%%%%%%%%%%%%%%%%%%%%%%%%%%%%%%%%%%

\begin{theorem}[Monotone Convergence Theorem]
Suppose

\[
0\leq X_1\leq X_2\leq\cdots
\]

and

\[
X_n\uparrow X
\]

almost surely.

Then

\[
\boxed{
\mathbb{E}[X_n]
\uparrow
\mathbb{E}[X].
}
\]
\end{theorem}

This allows limits and expectations to be interchanged for increasing
nonnegative sequences.

%%%%%%%%%%%%%%%%%%%%%%%%%%%%%%%%%%%%%%%%%%%%%%%%%%%%%%%%%%%%
\subsection{Worked Example: Truncating a Nonnegative Variable}
\label{subsec:worked-mct-truncation}
%%%%%%%%%%%%%%%%%%%%%%%%%%%%%%%%%%%%%%%%%%%%%%%%%%%%%%%%%%%%

\begin{workedexamplebox}{Truncated Random Variables}

Let

\[
X\geq0
\]

and define

\[
X_n
=
\min(X,n).
\]

Then

\[
X_n\uparrow X.
\]

By monotone convergence,

\begin{answerbox}
\[
\boxed{
\mathbb{E}[X_n]
\to
\mathbb{E}[X].
}
\]
\end{answerbox}

This remains valid even if

\[
\mathbb{E}[X]=\infty.
\]

\end{workedexamplebox}

%%%%%%%%%%%%%%%%%%%%%%%%%%%%%%%%%%%%%%%%%%%%%%%%%%%%%%%%%%%%
\subsection{Fatou's Lemma}
\label{subsec:fatou-probability}
%%%%%%%%%%%%%%%%%%%%%%%%%%%%%%%%%%%%%%%%%%%%%%%%%%%%%%%%%%%%

\begin{theorem}[Fatou's Lemma]
If

\[
X_n\geq0,
\]

then

\[
\boxed{
\mathbb{E}
\left[
\liminf_{n\to\infty}X_n
\right]
\leq
\liminf_{n\to\infty}
\mathbb{E}[X_n].
}
\]
\end{theorem}

The notation

\[
\liminf
\]

means ``limit inferior.''

Fatou's lemma provides a one-sided inequality when direct interchange of
limit and expectation is not justified.

%%%%%%%%%%%%%%%%%%%%%%%%%%%%%%%%%%%%%%%%%%%%%%%%%%%%%%%%%%%%
\subsection{Dominated Convergence Theorem}
\label{subsec:dct-probability}
%%%%%%%%%%%%%%%%%%%%%%%%%%%%%%%%%%%%%%%%%%%%%%%%%%%%%%%%%%%%

\begin{theorem}[Dominated Convergence Theorem]
Suppose

\[
X_n\to X
\]

almost surely and there exists an integrable random variable \(Y\) such
that

\[
|X_n|\leq Y
\]

almost surely for every \(n\).

Then

\[
\boxed{
\mathbb{E}[X_n]
\to
\mathbb{E}[X].
}
\]
\end{theorem}

Moreover,

\[
\boxed{
\mathbb{E}[|X_n-X|]\to0.
}
\]

%%%%%%%%%%%%%%%%%%%%%%%%%%%%%%%%%%%%%%%%%%%%%%%%%%%%%%%%%%%%
\subsection{Why Dominated Convergence Matters}
\label{subsec:why-dct-probability}
%%%%%%%%%%%%%%%%%%%%%%%%%%%%%%%%%%%%%%%%%%%%%%%%%%%%%%%%%%%%

A common but invalid step is

\[
\mathbb{E}
\left[
\lim_{n\to\infty}X_n
\right]
=
\lim_{n\to\infty}
\mathbb{E}[X_n]
\]

without justification.

Dominated convergence provides conditions under which this interchange
is valid.

Thus,

\[
\boxed{
\text{pointwise or almost-sure convergence alone is not enough}.
}
\]

%%%%%%%%%%%%%%%%%%%%%%%%%%%%%%%%%%%%%%%%%%%%%%%%%%%%%%%%%%%%
\subsection{Worked Example: Dominated Convergence}
\label{subsec:worked-dct}
%%%%%%%%%%%%%%%%%%%%%%%%%%%%%%%%%%%%%%%%%%%%%%%%%%%%%%%%%%%%

\begin{workedexamplebox}{A Bounded Sequence}

Suppose

\[
X_n
=
\frac{X}{n}
\]

and

\[
\mathbb{E}[|X|]<\infty.
\]

Then

\[
X_n\to0
\]

almost surely.

Also,

\[
|X_n|
\leq
|X|.
\]

Since \(|X|\) is integrable, dominated convergence gives

\[
\mathbb{E}[X_n]
\to0.
\]

\begin{answerbox}
\[
\boxed{
\lim_{n\to\infty}
\mathbb{E}
\left[
\frac{X}{n}
\right]
=
0.
}
\]
\end{answerbox}

\end{workedexamplebox}

%%%%%%%%%%%%%%%%%%%%%%%%%%%%%%%%%%%%%%%%%%%%%%%%%%%%%%%%%%%%
\subsection{Tonelli's Theorem in Probability}
\label{subsec:tonelli-probability}
%%%%%%%%%%%%%%%%%%%%%%%%%%%%%%%%%%%%%%%%%%%%%%%%%%%%%%%%%%%%

For nonnegative measurable random quantities, Tonelli's theorem allows
the order of integration to be exchanged.

If

\[
f(\omega,t)\geq0,
\]

then under the product-measure framework,

\[
\boxed{
\int
\left(
\int f(\omega,t)\,dt
\right)
dP(\omega)
=
\int
\left(
\int f(\omega,t)\,dP(\omega)
\right)
dt.
}
\]

This is frequently used in tail-integral formulas.

%%%%%%%%%%%%%%%%%%%%%%%%%%%%%%%%%%%%%%%%%%%%%%%%%%%%%%%%%%%%
\subsection{Tail Integral Formula Measure-Theoretically}
\label{subsec:tail-integral-measure}
%%%%%%%%%%%%%%%%%%%%%%%%%%%%%%%%%%%%%%%%%%%%%%%%%%%%%%%%%%%%

For

\[
X\geq0,
\]

write

\[
X(\omega)
=
\int_0^\infty
\mathbf{1}_{\{X(\omega)>t\}}
\,dt.
\]

Taking expectations and applying Tonelli,

\[
\mathbb{E}[X]
=
\int_0^\infty
P(X>t)\,dt.
\]

Thus,

\[
\boxed{
\mathbb{E}[X]
=
\int_0^\infty
P(X>t)\,dt.
}
\]

%%%%%%%%%%%%%%%%%%%%%%%%%%%%%%%%%%%%%%%%%%%%%%%%%%%%%%%%%%%%
\subsection{\(L^p\) Spaces of Random Variables}
\label{subsec:lp-probability}
%%%%%%%%%%%%%%%%%%%%%%%%%%%%%%%%%%%%%%%%%%%%%%%%%%%%%%%%%%%%

For

\[
1\leq p<\infty,
\]

define

\[
\boxed{
L^p(\Omega,\mathcal{F},P)
=
\left\{
X:
\mathbb{E}[|X|^p]<\infty
\right\}
}
\]

where random variables equal almost surely are identified.

The norm is

\[
\boxed{
\|X\|_p
=
\left(
\mathbb{E}[|X|^p]
\right)^{1/p}.
}
\]

%%%%%%%%%%%%%%%%%%%%%%%%%%%%%%%%%%%%%%%%%%%%%%%%%%%%%%%%%%%%
\subsection{\(L^1\) and \(L^2\)}
\label{subsec:l1-l2-probability}
%%%%%%%%%%%%%%%%%%%%%%%%%%%%%%%%%%%%%%%%%%%%%%%%%%%%%%%%%%%%

The space

\[
L^1
\]

contains integrable random variables:

\[
\boxed{
\mathbb{E}[|X|]<\infty.
}
\]

The space

\[
L^2
\]

contains square-integrable variables:

\[
\boxed{
\mathbb{E}[X^2]<\infty.
}
\]

Variance naturally belongs to the \(L^2\) setting.

%%%%%%%%%%%%%%%%%%%%%%%%%%%%%%%%%%%%%%%%%%%%%%%%%%%%%%%%%%%%
\subsection{\(L^p\) Convergence}
\label{subsec:lp-convergence-probability}
%%%%%%%%%%%%%%%%%%%%%%%%%%%%%%%%%%%%%%%%%%%%%%%%%%%%%%%%%%%%

\begin{definitionbox}{Convergence in \(L^p\)}
A sequence \(X_n\) converges to \(X\) in \(L^p\) if

\[
\boxed{
\mathbb{E}
[
|X_n-X|^p
]
\to0.
}
\]
\end{definitionbox}

We write

\[
\boxed{
X_n
\xrightarrow{L^p}
X.
}
\]

For \(p=2\), this is called mean-square convergence.

%%%%%%%%%%%%%%%%%%%%%%%%%%%%%%%%%%%%%%%%%%%%%%%%%%%%%%%%%%%%
\subsection{\(L^p\) Convergence Implies Convergence in Probability}
\label{subsec:lp-implies-probability}
%%%%%%%%%%%%%%%%%%%%%%%%%%%%%%%%%%%%%%%%%%%%%%%%%%%%%%%%%%%%

Suppose

\[
X_n
\xrightarrow{L^p}
X
\]

for some

\[
p>0.
\]

By Markov's inequality,

\[
P(|X_n-X|>\varepsilon)
=
P(|X_n-X|^p>\varepsilon^p).
\]

Therefore,

\[
\leq
\frac{
\mathbb{E}[|X_n-X|^p]
}{
\varepsilon^p
}.
\]

Hence,

\[
\boxed{
X_n
\xrightarrow{L^p}
X
\Longrightarrow
X_n
\xrightarrow{P}
X.
}
\]

%%%%%%%%%%%%%%%%%%%%%%%%%%%%%%%%%%%%%%%%%%%%%%%%%%%%%%%%%%%%
\subsection{Convergence Modes}
\label{subsec:convergence-modes-measure}
%%%%%%%%%%%%%%%%%%%%%%%%%%%%%%%%%%%%%%%%%%%%%%%%%%%%%%%%%%%%

Important modes include:

\[
\boxed{
X_n\xrightarrow{\text{a.s.}}X,
}
\]

\[
\boxed{
X_n\xrightarrow{P}X,
}
\]

\[
\boxed{
X_n\xrightarrow{d}X,
}
\]

and

\[
\boxed{
X_n\xrightarrow{L^p}X.
}
\]

Some useful implications are

\[
\boxed{
X_n\xrightarrow{\text{a.s.}}X
\Longrightarrow
X_n\xrightarrow{P}X,
}
\]

and

\[
\boxed{
X_n\xrightarrow{L^p}X
\Longrightarrow
X_n\xrightarrow{P}X.
}
\]

Also,

\[
\boxed{
X_n\xrightarrow{P}X
\Longrightarrow
X_n\xrightarrow{d}X.
}
\]

%%%%%%%%%%%%%%%%%%%%%%%%%%%%%%%%%%%%%%%%%%%%%%%%%%%%%%%%%%%%
\subsection{No General Implication Between a.s. and \(L^p\)}
\label{subsec:as-lp-no-general}
%%%%%%%%%%%%%%%%%%%%%%%%%%%%%%%%%%%%%%%%%%%%%%%%%%%%%%%%%%%%

Almost-sure convergence does not by itself imply \(L^p\) convergence.

Similarly, \(L^p\) convergence does not by itself imply convergence
almost surely along the full sequence.

Additional assumptions are required.

For example, dominated convergence can convert almost-sure convergence
into \(L^1\) convergence when a common integrable bound exists.

%%%%%%%%%%%%%%%%%%%%%%%%%%%%%%%%%%%%%%%%%%%%%%%%%%%%%%%%%%%%
\subsection{Subsequence Principle}
\label{subsec:subsequence-principle}
%%%%%%%%%%%%%%%%%%%%%%%%%%%%%%%%%%%%%%%%%%%%%%%%%%%%%%%%%%%%

If

\[
X_n\xrightarrow{P}X,
\]

then every subsequence contains a further subsequence that converges to
\(X\) almost surely.

This gives an important connection between convergence in probability
and almost-sure convergence.

%%%%%%%%%%%%%%%%%%%%%%%%%%%%%%%%%%%%%%%%%%%%%%%%%%%%%%%%%%%%
\subsection{Sigma-Algebra Generated by a Random Variable}
\label{subsec:sigma-generated-rv}
%%%%%%%%%%%%%%%%%%%%%%%%%%%%%%%%%%%%%%%%%%%%%%%%%%%%%%%%%%%%

A random variable \(X\) generates a sigma-algebra

\[
\boxed{
\sigma(X).
}
\]

This is the smallest sigma-algebra that makes \(X\) measurable.

It represents all information that can be determined from knowing \(X\).

For example,

\[
\{X\leq x\}
\in
\sigma(X).
\]

%%%%%%%%%%%%%%%%%%%%%%%%%%%%%%%%%%%%%%%%%%%%%%%%%%%%%%%%%%%%
\subsection{Information Interpretation}
\label{subsec:sigma-algebra-information}
%%%%%%%%%%%%%%%%%%%%%%%%%%%%%%%%%%%%%%%%%%%%%%%%%%%%%%%%%%%%

Sigma-algebras can be interpreted as information sets.

A larger sigma-algebra contains more events and therefore represents
more information.

If

\[
\mathcal{G}
\subseteq
\mathcal{F},
\]

then \(\mathcal{G}\) represents partial information relative to the full
information \(\mathcal{F}\).

This viewpoint is essential for conditional expectation and stochastic
processes.

%%%%%%%%%%%%%%%%%%%%%%%%%%%%%%%%%%%%%%%%%%%%%%%%%%%%%%%%%%%%
\subsection{Conditional Expectation with Respect to a Sigma-Algebra}
\label{subsec:conditional-expectation-sigma}
%%%%%%%%%%%%%%%%%%%%%%%%%%%%%%%%%%%%%%%%%%%%%%%%%%%%%%%%%%%%

\begin{definitionbox}{Conditional Expectation}
Let

\[
X\in L^1(\Omega,\mathcal{F},P),
\]

and let

\[
\mathcal{G}
\subseteq
\mathcal{F}
\]

be a sub-sigma-algebra.

A conditional expectation of \(X\) given \(\mathcal{G}\) is a
\(\mathcal{G}\)-measurable random variable

\[
\boxed{
\mathbb{E}[X\mid\mathcal{G}]
}
\]

such that

\[
\boxed{
\int_G
\mathbb{E}[X\mid\mathcal{G}]
\,dP
=
\int_G
X\,dP
}
\]

for every

\[
G\in\mathcal{G}.
\]
\end{definitionbox}

%%%%%%%%%%%%%%%%%%%%%%%%%%%%%%%%%%%%%%%%%%%%%%%%%%%%%%%%%%%%
\subsection{Interpretation of Conditional Expectation}
\label{subsec:conditional-expectation-information}
%%%%%%%%%%%%%%%%%%%%%%%%%%%%%%%%%%%%%%%%%%%%%%%%%%%%%%%%%%%%

The quantity

\[
\mathbb{E}[X\mid\mathcal{G}]
\]

is the best expectation of \(X\) using only the information contained in
\(\mathcal{G}\).

It must be measurable with respect to \(\mathcal{G}\), meaning it cannot
use information unavailable in \(\mathcal{G}\).

%%%%%%%%%%%%%%%%%%%%%%%%%%%%%%%%%%%%%%%%%%%%%%%%%%%%%%%%%%%%
\subsection{Conditional Expectation Given a Random Variable}
\label{subsec:conditional-expectation-given-rv}
%%%%%%%%%%%%%%%%%%%%%%%%%%%%%%%%%%%%%%%%%%%%%%%%%%%%%%%%%%%%

Conditioning on a random variable \(Y\) means conditioning on the
sigma-algebra generated by \(Y\):

\[
\boxed{
\mathbb{E}[X\mid Y]
=
\mathbb{E}[X\mid\sigma(Y)].
}
\]

Thus the familiar notation

\[
\mathbb{E}[X\mid Y]
\]

is shorthand for a sigma-algebra conditional expectation.

%%%%%%%%%%%%%%%%%%%%%%%%%%%%%%%%%%%%%%%%%%%%%%%%%%%%%%%%%%%%
\subsection{Conditional Probability}
\label{subsec:conditional-probability-measure}
%%%%%%%%%%%%%%%%%%%%%%%%%%%%%%%%%%%%%%%%%%%%%%%%%%%%%%%%%%%%

For an event \(A\), define

\[
\boxed{
P(A\mid\mathcal{G})
=
\mathbb{E}
[
\mathbf{1}_A
\mid
\mathcal{G}
].
}
\]

Thus conditional probability is a special case of conditional
expectation.

This unifies elementary conditioning with modern probability theory.

%%%%%%%%%%%%%%%%%%%%%%%%%%%%%%%%%%%%%%%%%%%%%%%%%%%%%%%%%%%%
\subsection{Basic Properties of Conditional Expectation}
\label{subsec:conditional-expectation-properties}
%%%%%%%%%%%%%%%%%%%%%%%%%%%%%%%%%%%%%%%%%%%%%%%%%%%%%%%%%%%%

Conditional expectation satisfies linearity:

\[
\boxed{
\mathbb{E}[aX+bY\mid\mathcal{G}]
=
a\mathbb{E}[X\mid\mathcal{G}]
+
b\mathbb{E}[Y\mid\mathcal{G}].
}
\]

If \(X\geq0\), then

\[
\boxed{
\mathbb{E}[X\mid\mathcal{G}]
\geq0
}
\]

almost surely.

If \(X\) is already \(\mathcal{G}\)-measurable, then

\[
\boxed{
\mathbb{E}[X\mid\mathcal{G}]
=
X
}
\]

almost surely.

%%%%%%%%%%%%%%%%%%%%%%%%%%%%%%%%%%%%%%%%%%%%%%%%%%%%%%%%%%%%
\subsection{Pulling Out Known Factors}
\label{subsec:pull-out-known-factor}
%%%%%%%%%%%%%%%%%%%%%%%%%%%%%%%%%%%%%%%%%%%%%%%%%%%%%%%%%%%%

If \(Y\) is \(\mathcal{G}\)-measurable and the relevant quantities are
integrable, then

\[
\boxed{
\mathbb{E}[XY\mid\mathcal{G}]
=
Y
\mathbb{E}[X\mid\mathcal{G}].
}
\]

The variable \(Y\) is ``known'' given \(\mathcal{G}\), so it behaves like
a constant inside the conditional expectation.

%%%%%%%%%%%%%%%%%%%%%%%%%%%%%%%%%%%%%%%%%%%%%%%%%%%%%%%%%%%%
\subsection{Tower Property}
\label{subsec:tower-property-measure}
%%%%%%%%%%%%%%%%%%%%%%%%%%%%%%%%%%%%%%%%%%%%%%%%%%%%%%%%%%%%

Suppose

\[
\mathcal{H}
\subseteq
\mathcal{G}
\subseteq
\mathcal{F}.
\]

Then

\[
\boxed{
\mathbb{E}
[
\mathbb{E}[X\mid\mathcal{G}]
\mid
\mathcal{H}
]
=
\mathbb{E}[X\mid\mathcal{H}].
}
\]

In particular,

\[
\boxed{
\mathbb{E}
[
\mathbb{E}[X\mid\mathcal{G}]
]
=
\mathbb{E}[X].
}
\]

This is the measure-theoretic law of total expectation.

%%%%%%%%%%%%%%%%%%%%%%%%%%%%%%%%%%%%%%%%%%%%%%%%%%%%%%%%%%%%
\subsection{Worked Example: Conditioning on Trivial Information}
\label{subsec:worked-trivial-conditioning}
%%%%%%%%%%%%%%%%%%%%%%%%%%%%%%%%%%%%%%%%%%%%%%%%%%%%%%%%%%%%

\begin{workedexamplebox}{The Trivial Sigma-Algebra}

Let

\[
\mathcal{G}
=
\{\varnothing,\Omega\}.
\]

This sigma-algebra contains no nontrivial information.

Any \(\mathcal{G}\)-measurable random variable must be constant almost
surely.

Therefore,

\begin{answerbox}
\[
\boxed{
\mathbb{E}[X\mid\mathcal{G}]
=
\mathbb{E}[X].
}
\]
\end{answerbox}

\end{workedexamplebox}

%%%%%%%%%%%%%%%%%%%%%%%%%%%%%%%%%%%%%%%%%%%%%%%%%%%%%%%%%%%%
\subsection{Conditioning on Full Information}
\label{subsec:conditioning-full-information}
%%%%%%%%%%%%%%%%%%%%%%%%%%%%%%%%%%%%%%%%%%%%%%%%%%%%%%%%%%%%

If

\[
\mathcal{G}
=
\mathcal{F},
\]

then \(X\) itself is \(\mathcal{G}\)-measurable.

Therefore,

\[
\boxed{
\mathbb{E}[X\mid\mathcal{F}]
=
X
}
\]

almost surely.

Thus:

\[
\boxed{
\text{no information}
\longrightarrow
\mathbb{E}[X],
}
\]

while

\[
\boxed{
\text{complete information}
\longrightarrow
X.
}
\]

%%%%%%%%%%%%%%%%%%%%%%%%%%%%%%%%%%%%%%%%%%%%%%%%%%%%%%%%%%%%
\subsection{Conditional Expectation as Projection}
\label{subsec:conditional-expectation-projection}
%%%%%%%%%%%%%%%%%%%%%%%%%%%%%%%%%%%%%%%%%%%%%%%%%%%%%%%%%%%%

Suppose

\[
X\in L^2.
\]

Then

\[
\mathbb{E}[X\mid\mathcal{G}]
\]

is the orthogonal projection of \(X\) onto the closed subspace of
\(\mathcal{G}\)-measurable square-integrable random variables.

Thus,

\[
\boxed{
\mathbb{E}
\left[
(X-\mathbb{E}[X\mid\mathcal{G}])Y
\right]
=
0
}
\]

for every square-integrable \(\mathcal{G}\)-measurable \(Y\).

%%%%%%%%%%%%%%%%%%%%%%%%%%%%%%%%%%%%%%%%%%%%%%%%%%%%%%%%%%%%
\subsection{Best Mean-Square Predictor}
\label{subsec:best-mean-square-predictor}
%%%%%%%%%%%%%%%%%%%%%%%%%%%%%%%%%%%%%%%%%%%%%%%%%%%%%%%%%%%%

Among all \(\mathcal{G}\)-measurable square-integrable random variables
\(Y\),

\[
\boxed{
\mathbb{E}[X\mid\mathcal{G}]
}
\]

minimizes

\[
\mathbb{E}[(X-Y)^2].
\]

Thus conditional expectation is the optimal predictor of \(X\) using
information \(\mathcal{G}\) under squared-error loss.

%%%%%%%%%%%%%%%%%%%%%%%%%%%%%%%%%%%%%%%%%%%%%%%%%%%%%%%%%%%%
\subsection{Law of Total Variance Measure-Theoretically}
\label{subsec:total-variance-measure}
%%%%%%%%%%%%%%%%%%%%%%%%%%%%%%%%%%%%%%%%%%%%%%%%%%%%%%%%%%%%

For square-integrable \(X\),

\[
\boxed{
\operatorname{Var}(X)
=
\mathbb{E}
[
\operatorname{Var}(X\mid\mathcal{G})
]
+
\operatorname{Var}
(
\mathbb{E}[X\mid\mathcal{G}]
).
}
\]

This is the orthogonal decomposition of total variation into:

\[
\boxed{
\text{unexplained conditional variation}
}
\]

and

\[
\boxed{
\text{variation explained by }\mathcal{G}.
}
\]

%%%%%%%%%%%%%%%%%%%%%%%%%%%%%%%%%%%%%%%%%%%%%%%%%%%%%%%%%%%%
\subsection{Radon--Nikodym Theorem}
\label{subsec:radon-nikodym-probability}
%%%%%%%%%%%%%%%%%%%%%%%%%%%%%%%%%%%%%%%%%%%%%%%%%%%%%%%%%%%%

A central theorem underlying conditional expectation is the
Radon--Nikodym theorem.

\begin{theorem}[Radon--Nikodym Theorem]
Suppose \(\nu\) and \(\mu\) are sigma-finite measures and

\[
\nu\ll\mu,
\]

meaning \(\nu\) is absolutely continuous with respect to \(\mu\).

Then there exists a measurable function \(f\) such that

\[
\boxed{
\nu(A)
=
\int_A
f\,d\mu
}
\]

for measurable \(A\).

The function is written

\[
\boxed{
f
=
\frac{d\nu}{d\mu}.
}
\]
\end{theorem}

%%%%%%%%%%%%%%%%%%%%%%%%%%%%%%%%%%%%%%%%%%%%%%%%%%%%%%%%%%%%
\subsection{Conditional Expectation via Radon--Nikodym}
\label{subsec:conditional-radon-nikodym}
%%%%%%%%%%%%%%%%%%%%%%%%%%%%%%%%%%%%%%%%%%%%%%%%%%%%%%%%%%%%

For integrable \(X\geq0\), define a measure on \(\mathcal{G}\) by

\[
\nu(G)
=
\int_G
X\,dP.
\]

Because

\[
\nu\ll P|_{\mathcal{G}},
\]

the Radon--Nikodym theorem gives a \(\mathcal{G}\)-measurable function
\(Y\) such that

\[
\nu(G)
=
\int_G
Y\,dP.
\]

That function is

\[
\boxed{
Y
=
\mathbb{E}[X\mid\mathcal{G}].
}
\]

Thus conditional expectation exists through a measure-theoretic
derivative construction.

%%%%%%%%%%%%%%%%%%%%%%%%%%%%%%%%%%%%%%%%%%%%%%%%%%%%%%%%%%%%
\subsection{Independence of Sigma-Algebras}
\label{subsec:independence-sigma-algebras}
%%%%%%%%%%%%%%%%%%%%%%%%%%%%%%%%%%%%%%%%%%%%%%%%%%%%%%%%%%%%

\begin{definitionbox}{Independent Sigma-Algebras}
Sigma-algebras

\[
\mathcal{F}_1,\ldots,\mathcal{F}_n
\]

are independent if for every choice

\[
A_i\in\mathcal{F}_i,
\]

\[
\boxed{
P
\left(
\bigcap_{i=1}^{n}A_i
\right)
=
\prod_{i=1}^{n}P(A_i).
}
\]
\end{definitionbox}

This generalizes independence from events to entire information
structures.

%%%%%%%%%%%%%%%%%%%%%%%%%%%%%%%%%%%%%%%%%%%%%%%%%%%%%%%%%%%%
\subsection{Independence of Random Variables}
\label{subsec:rv-independence-measure}
%%%%%%%%%%%%%%%%%%%%%%%%%%%%%%%%%%%%%%%%%%%%%%%%%%%%%%%%%%%%

Random variables

\[
X_1,\ldots,X_n
\]

are independent if the sigma-algebras

\[
\sigma(X_1),\ldots,\sigma(X_n)
\]

are independent.

Equivalently, for Borel sets

\[
B_1,\ldots,B_n,
\]

\[
\boxed{
P
\left(
X_1\in B_1,\ldots,X_n\in B_n
\right)
=
\prod_{i=1}^{n}
P(X_i\in B_i).
}
\]

%%%%%%%%%%%%%%%%%%%%%%%%%%%%%%%%%%%%%%%%%%%%%%%%%%%%%%%%%%%%
\subsection{Independence and Conditional Expectation}
\label{subsec:independence-conditional-expectation}
%%%%%%%%%%%%%%%%%%%%%%%%%%%%%%%%%%%%%%%%%%%%%%%%%%%%%%%%%%%%

If \(X\) is integrable and independent of \(\mathcal{G}\), then

\[
\boxed{
\mathbb{E}[X\mid\mathcal{G}]
=
\mathbb{E}[X]
}
\]

almost surely.

Thus conditioning on independent information provides no change in the
best prediction of \(X\).

%%%%%%%%%%%%%%%%%%%%%%%%%%%%%%%%%%%%%%%%%%%%%%%%%%%%%%%%%%%%
\subsection{Product Probability Spaces}
\label{subsec:product-probability-spaces}
%%%%%%%%%%%%%%%%%%%%%%%%%%%%%%%%%%%%%%%%%%%%%%%%%%%%%%%%%%%%

Suppose

\[
(\Omega_1,\mathcal{F}_1,P_1)
\]

and

\[
(\Omega_2,\mathcal{F}_2,P_2)
\]

are probability spaces.

Their product space is

\[
\boxed{
(\Omega_1\times\Omega_2,
\mathcal{F}_1\otimes\mathcal{F}_2,
P_1\otimes P_2).
}
\]

This framework constructs independent coordinates naturally.

%%%%%%%%%%%%%%%%%%%%%%%%%%%%%%%%%%%%%%%%%%%%%%%%%%%%%%%%%%%%
\subsection{Coordinate Random Variables on Product Spaces}
\label{subsec:coordinate-rvs-product}
%%%%%%%%%%%%%%%%%%%%%%%%%%%%%%%%%%%%%%%%%%%%%%%%%%%%%%%%%%%%

On

\[
\Omega_1\times\Omega_2,
\]

define

\[
X(\omega_1,\omega_2)
=
\omega_1
\]

and

\[
Y(\omega_1,\omega_2)
=
\omega_2.
\]

Under the product measure, the coordinate variables are independent.

Thus product spaces provide a canonical construction of independent
random variables.

%%%%%%%%%%%%%%%%%%%%%%%%%%%%%%%%%%%%%%%%%%%%%%%%%%%%%%%%%%%%
\subsection{Fubini's Theorem in Probability}
\label{subsec:fubini-probability}
%%%%%%%%%%%%%%%%%%%%%%%%%%%%%%%%%%%%%%%%%%%%%%%%%%%%%%%%%%%%

If a jointly measurable function \(g(X,Y)\) is integrable under a
product probability measure, Fubini's theorem allows

\[
\boxed{
\mathbb{E}[g(X,Y)]
=
\int
\left(
\int
g(x,y)\,dP_Y(y)
\right)
dP_X(x).
}
\]

The order of integration may be reversed.

This is foundational for joint-distribution calculations.

%%%%%%%%%%%%%%%%%%%%%%%%%%%%%%%%%%%%%%%%%%%%%%%%%%%%%%%%%%%%
\subsection{Infinite Product Spaces Preview}
\label{subsec:infinite-product-spaces-preview}
%%%%%%%%%%%%%%%%%%%%%%%%%%%%%%%%%%%%%%%%%%%%%%%%%%%%%%%%%%%%

Sequences of random variables may be constructed on infinite product
spaces such as

\[
\Omega
=
\mathbb{R}^{\mathbb{N}}.
\]

An outcome is then an infinite sequence

\[
\omega
=
(\omega_1,\omega_2,\ldots).
\]

Coordinate maps

\[
X_n(\omega)
=
\omega_n
\]

represent the random variables.

This framework is fundamental for stochastic processes.

%%%%%%%%%%%%%%%%%%%%%%%%%%%%%%%%%%%%%%%%%%%%%%%%%%%%%%%%%%%%
\subsection{Borel--Cantelli Lemmas Revisited}
\label{subsec:borel-cantelli-measure}
%%%%%%%%%%%%%%%%%%%%%%%%%%%%%%%%%%%%%%%%%%%%%%%%%%%%%%%%%%%%

For events

\[
A_1,A_2,\ldots,
\]

define

\[
\boxed{
\limsup_{n\to\infty}A_n
=
\bigcap_{m=1}^{\infty}
\bigcup_{n=m}^{\infty}A_n.
}
\]

This is the event that infinitely many \(A_n\) occur.

The notation

\[
A_n\text{ i.o.}
\]

means the same thing.

%%%%%%%%%%%%%%%%%%%%%%%%%%%%%%%%%%%%%%%%%%%%%%%%%%%%%%%%%%%%
\subsection{First Borel--Cantelli Lemma}
\label{subsec:first-borel-cantelli-measure}
%%%%%%%%%%%%%%%%%%%%%%%%%%%%%%%%%%%%%%%%%%%%%%%%%%%%%%%%%%%%

If

\[
\sum_{n=1}^{\infty}
P(A_n)
<
\infty,
\]

then

\[
\boxed{
P(A_n\text{ i.o.})=0.
}
\]

No independence assumption is needed.

%%%%%%%%%%%%%%%%%%%%%%%%%%%%%%%%%%%%%%%%%%%%%%%%%%%%%%%%%%%%
\subsection{Second Borel--Cantelli Lemma}
\label{subsec:second-borel-cantelli-measure}
%%%%%%%%%%%%%%%%%%%%%%%%%%%%%%%%%%%%%%%%%%%%%%%%%%%%%%%%%%%%

If the events \(A_n\) are independent and

\[
\sum_{n=1}^{\infty}
P(A_n)
=
\infty,
\]

then

\[
\boxed{
P(A_n\text{ i.o.})=1.
}
\]

The independence assumption is crucial here.

%%%%%%%%%%%%%%%%%%%%%%%%%%%%%%%%%%%%%%%%%%%%%%%%%%%%%%%%%%%%
\subsection{Worked Example: Rare Events}
\label{subsec:worked-borel-cantelli}
%%%%%%%%%%%%%%%%%%%%%%%%%%%%%%%%%%%%%%%%%%%%%%%%%%%%%%%%%%%%

\begin{workedexamplebox}{Events with Summable Probabilities}

Suppose

\[
P(A_n)
=
\frac1{n^2}.
\]

Since

\[
\sum_{n=1}^{\infty}
\frac1{n^2}
<
\infty,
\]

the first Borel--Cantelli lemma gives

\begin{answerbox}
\[
\boxed{
P(A_n\text{ i.o.})=0.
}
\]
\end{answerbox}

Thus only finitely many of the \(A_n\) occur almost surely.

\end{workedexamplebox}

%%%%%%%%%%%%%%%%%%%%%%%%%%%%%%%%%%%%%%%%%%%%%%%%%%%%%%%%%%%%
\subsection{Regular Conditional Probability Preview}
\label{subsec:regular-conditional-probability}
%%%%%%%%%%%%%%%%%%%%%%%%%%%%%%%%%%%%%%%%%%%%%%%%%%%%%%%%%%%%

Elementary notation such as

\[
P(X\in A\mid Y=y)
\]

can be delicate when

\[
P(Y=y)=0.
\]

A regular conditional probability provides a probability measure
depending measurably on the conditioning value.

Informally,

\[
\boxed{
y
\mapsto
P(X\in A\mid Y=y)
}
\]

behaves like a conditional distribution for each \(y\).

Existence requires suitable measurable-space assumptions.

%%%%%%%%%%%%%%%%%%%%%%%%%%%%%%%%%%%%%%%%%%%%%%%%%%%%%%%%%%%%
\subsection{Probability Kernels Preview}
\label{subsec:probability-kernels-preview}
%%%%%%%%%%%%%%%%%%%%%%%%%%%%%%%%%%%%%%%%%%%%%%%%%%%%%%%%%%%%

A probability kernel assigns a probability measure to each point of
another measurable space.

One may write

\[
\boxed{
K(y,A)
=
P(X\in A\mid Y=y).
}
\]

For fixed \(y\),

\[
A\mapsto K(y,A)
\]

is a probability measure.

For fixed \(A\),

\[
y\mapsto K(y,A)
\]

is measurable.

Kernels are fundamental in Markov processes and Bayesian probability.

%%%%%%%%%%%%%%%%%%%%%%%%%%%%%%%%%%%%%%%%%%%%%%%%%%%%%%%%%%%%
\subsection{Filtrations Preview}
\label{subsec:filtrations-preview}
%%%%%%%%%%%%%%%%%%%%%%%%%%%%%%%%%%%%%%%%%%%%%%%%%%%%%%%%%%%%

A filtration is an increasing family of sigma-algebras:

\[
\boxed{
\mathcal{F}_0
\subseteq
\mathcal{F}_1
\subseteq
\mathcal{F}_2
\subseteq
\cdots.
}
\]

The interpretation is that

\[
\mathcal{F}_n
\]

contains the information available by time \(n\).

Filtrations are central to martingales, stopping times, and stochastic
calculus.

%%%%%%%%%%%%%%%%%%%%%%%%%%%%%%%%%%%%%%%%%%%%%%%%%%%%%%%%%%%%
\subsection{Adapted Random Processes Preview}
\label{subsec:adapted-processes-preview}
%%%%%%%%%%%%%%%%%%%%%%%%%%%%%%%%%%%%%%%%%%%%%%%%%%%%%%%%%%%%

A stochastic process

\[
X_0,X_1,X_2,\ldots
\]

is adapted to a filtration

\[
(\mathcal{F}_n)
\]

if

\[
\boxed{
X_n
\text{ is }\mathcal{F}_n\text{-measurable}
}
\]

for every \(n\).

Thus the value \(X_n\) is determined by information available at time
\(n\).

%%%%%%%%%%%%%%%%%%%%%%%%%%%%%%%%%%%%%%%%%%%%%%%%%%%%%%%%%%%%
\subsection{Martingales Preview}
\label{subsec:martingales-measure-preview}
%%%%%%%%%%%%%%%%%%%%%%%%%%%%%%%%%%%%%%%%%%%%%%%%%%%%%%%%%%%%

A process

\[
(M_n)
\]

adapted to \((\mathcal{F}_n)\) is a martingale if

\[
\mathbb{E}[|M_n|]<\infty
\]

and

\[
\boxed{
\mathbb{E}[M_{n+1}\mid\mathcal{F}_n]
=
M_n.
}
\]

This formalizes the idea of a fair game.

Martingale theory depends fundamentally on measure-theoretic conditional
expectation.

%%%%%%%%%%%%%%%%%%%%%%%%%%%%%%%%%%%%%%%%%%%%%%%%%%%%%%%%%%%%
\subsection{Probability Spaces with Completion}
\label{subsec:completed-probability-space}
%%%%%%%%%%%%%%%%%%%%%%%%%%%%%%%%%%%%%%%%%%%%%%%%%%%%%%%%%%%%

A probability space is complete if every subset of every probability-zero
event is measurable.

That is, if

\[
A\in\mathcal{F},
\qquad
P(A)=0,
\]

and

\[
B\subseteq A,
\]

then completeness requires

\[
\boxed{
B\in\mathcal{F}.
}
\]

Probability spaces are often completed to avoid pathological
measurability issues.

%%%%%%%%%%%%%%%%%%%%%%%%%%%%%%%%%%%%%%%%%%%%%%%%%%%%%%%%%%%%
\subsection{Null Sets}
\label{subsec:null-sets-probability}
%%%%%%%%%%%%%%%%%%%%%%%%%%%%%%%%%%%%%%%%%%%%%%%%%%%%%%%%%%%%

A measurable event \(N\) is a null set if

\[
\boxed{
P(N)=0.
}
\]

A property holding outside a null set holds almost surely.

Null events need not be impossible in a pointwise sense.

For example, if \(X\) is continuously distributed, then

\[
P(X=x)=0
\]

for each fixed \(x\), even though \(X\) must take some value.

%%%%%%%%%%%%%%%%%%%%%%%%%%%%%%%%%%%%%%%%%%%%%%%%%%%%%%%%%%%%
\subsection{Essential Supremum Preview}
\label{subsec:essential-supremum-preview}
%%%%%%%%%%%%%%%%%%%%%%%%%%%%%%%%%%%%%%%%%%%%%%%%%%%%%%%%%%%%

The \(L^\infty\) norm is defined using the essential supremum:

\[
\boxed{
\|X\|_\infty
=
\operatorname*{ess\,sup}|X|.
}
\]

This ignores behavior on probability-zero sets.

Thus two random variables that differ only on a null set have the same
\(L^\infty\) norm.

%%%%%%%%%%%%%%%%%%%%%%%%%%%%%%%%%%%%%%%%%%%%%%%%%%%%%%%%%%%%
\subsection{Jensen's Inequality Measure-Theoretically}
\label{subsec:jensen-measure-probability}
%%%%%%%%%%%%%%%%%%%%%%%%%%%%%%%%%%%%%%%%%%%%%%%%%%%%%%%%%%%%

If \(\varphi\) is convex and \(X\) is integrable with appropriate
integrability of \(\varphi(X)\), then

\[
\boxed{
\varphi(\mathbb{E}[X])
\leq
\mathbb{E}[\varphi(X)].
}
\]

A conditional version also holds:

\[
\boxed{
\varphi
\left(
\mathbb{E}[X\mid\mathcal{G}]
\right)
\leq
\mathbb{E}
[
\varphi(X)
\mid
\mathcal{G}
]
}
\]

almost surely.

%%%%%%%%%%%%%%%%%%%%%%%%%%%%%%%%%%%%%%%%%%%%%%%%%%%%%%%%%%%%
\subsection{Hölder's Inequality}
\label{subsec:holder-probability}
%%%%%%%%%%%%%%%%%%%%%%%%%%%%%%%%%%%%%%%%%%%%%%%%%%%%%%%%%%%%

Suppose

\[
p,q>1
\]

satisfy

\[
\frac1p+\frac1q=1.
\]

Then

\[
\boxed{
\mathbb{E}[|XY|]
\leq
\|X\|_p
\|Y\|_q.
}
\]

This is Hölder's inequality in probability language.

%%%%%%%%%%%%%%%%%%%%%%%%%%%%%%%%%%%%%%%%%%%%%%%%%%%%%%%%%%%%
\subsection{Cauchy--Schwarz as a Special Case}
\label{subsec:cauchy-schwarz-measure-probability}
%%%%%%%%%%%%%%%%%%%%%%%%%%%%%%%%%%%%%%%%%%%%%%%%%%%%%%%%%%%%

Set

\[
p=q=2.
\]

Then Hölder's inequality gives

\[
\boxed{
|\mathbb{E}[XY]|
\leq
\sqrt{
\mathbb{E}[X^2]
\mathbb{E}[Y^2]
}.
}
\]

This is the Cauchy--Schwarz inequality for random variables.

%%%%%%%%%%%%%%%%%%%%%%%%%%%%%%%%%%%%%%%%%%%%%%%%%%%%%%%%%%%%
\subsection{Minkowski's Inequality}
\label{subsec:minkowski-probability}
%%%%%%%%%%%%%%%%%%%%%%%%%%%%%%%%%%%%%%%%%%%%%%%%%%%%%%%%%%%%

For

\[
p\geq1,
\]

Minkowski's inequality states

\[
\boxed{
\|X+Y\|_p
\leq
\|X\|_p+\|Y\|_p.
}
\]

Thus \(\|\cdot\|_p\) is genuinely a norm on \(L^p\).

This connects probability with functional analysis.

%%%%%%%%%%%%%%%%%%%%%%%%%%%%%%%%%%%%%%%%%%%%%%%%%%%%%%%%%%%%
\subsection{Uniform Integrability Preview}
\label{subsec:uniform-integrability-preview}
%%%%%%%%%%%%%%%%%%%%%%%%%%%%%%%%%%%%%%%%%%%%%%%%%%%%%%%%%%%%

A family of integrable random variables is uniformly integrable if its
expectations cannot hide significant mass in increasingly large tails.

One common condition is

\[
\boxed{
\lim_{K\to\infty}
\sup_n
\mathbb{E}
\left[
|X_n|
\mathbf{1}_{\{|X_n|>K\}}
\right]
=
0.
}
\]

Uniform integrability is crucial for upgrading convergence in
probability to convergence of expectations.

%%%%%%%%%%%%%%%%%%%%%%%%%%%%%%%%%%%%%%%%%%%%%%%%%%%%%%%%%%%%
\subsection{Vitali Convergence Preview}
\label{subsec:vitali-convergence-preview}
%%%%%%%%%%%%%%%%%%%%%%%%%%%%%%%%%%%%%%%%%%%%%%%%%%%%%%%%%%%%

A standard result states that if

\[
X_n\xrightarrow{P}X
\]

and \((X_n)\) is uniformly integrable, then under suitable conditions,

\[
\boxed{
\mathbb{E}[|X_n-X|]
\to0.
}
\]

Thus uniform integrability plays a role similar to domination but is
more flexible.

%%%%%%%%%%%%%%%%%%%%%%%%%%%%%%%%%%%%%%%%%%%%%%%%%%%%%%%%%%%%
\subsection{Tightness Preview}
\label{subsec:tightness-probability-preview}
%%%%%%%%%%%%%%%%%%%%%%%%%%%%%%%%%%%%%%%%%%%%%%%%%%%%%%%%%%%%

A family of probability measures

\[
\{\mu_\alpha\}
\]

on \(\mathbb{R}\) is tight if for every

\[
\varepsilon>0,
\]

there exists a compact set \(K\) such that

\[
\boxed{
\mu_\alpha(K)
\geq
1-\varepsilon
}
\]

for all \(\alpha\).

Tightness prevents probability mass from escaping to infinity.

It is fundamental in weak convergence theory.

%%%%%%%%%%%%%%%%%%%%%%%%%%%%%%%%%%%%%%%%%%%%%%%%%%%%%%%%%%%%
\subsection{Weak Convergence of Probability Measures}
\label{subsec:weak-convergence-measures}
%%%%%%%%%%%%%%%%%%%%%%%%%%%%%%%%%%%%%%%%%%%%%%%%%%%%%%%%%%%%

Convergence in distribution can be expressed as weak convergence of
probability measures.

If

\[
X_n\xrightarrow{d}X,
\]

then one writes

\[
\boxed{
P_{X_n}
\Rightarrow
P_X.
}
\]

A key characterization is that

\[
\boxed{
\int f\,dP_{X_n}
\to
\int f\,dP_X
}
\]

for every bounded continuous function \(f\).

%%%%%%%%%%%%%%%%%%%%%%%%%%%%%%%%%%%%%%%%%%%%%%%%%%%%%%%%%%%%
\subsection{Portmanteau Theorem Preview}
\label{subsec:portmanteau-preview}
%%%%%%%%%%%%%%%%%%%%%%%%%%%%%%%%%%%%%%%%%%%%%%%%%%%%%%%%%%%%

The Portmanteau theorem gives several equivalent characterizations of
weak convergence.

Among them:

\[
P_{X_n}\Rightarrow P_X
\]

if and only if

\[
\boxed{
\mathbb{E}[f(X_n)]
\to
\mathbb{E}[f(X)]
}
\]

for every bounded continuous \(f\).

It also gives inequalities involving closed and open sets.

%%%%%%%%%%%%%%%%%%%%%%%%%%%%%%%%%%%%%%%%%%%%%%%%%%%%%%%%%%%%
\subsection{Measure-Theoretic View of the LLN}
\label{subsec:measure-view-lln}
%%%%%%%%%%%%%%%%%%%%%%%%%%%%%%%%%%%%%%%%%%%%%%%%%%%%%%%%%%%%

The strong law of large numbers is an almost-sure statement:

\[
\boxed{
P
\left(
\left\{
\omega:
\frac1n
\sum_{i=1}^{n}
X_i(\omega)
\to
\mu
\right\}
\right)
=
1.
}
\]

Thus the strong law says a particular convergence set has full
probability measure.

%%%%%%%%%%%%%%%%%%%%%%%%%%%%%%%%%%%%%%%%%%%%%%%%%%%%%%%%%%%%
\subsection{Measure-Theoretic View of the CLT}
\label{subsec:measure-view-clt}
%%%%%%%%%%%%%%%%%%%%%%%%%%%%%%%%%%%%%%%%%%%%%%%%%%%%%%%%%%%%

The central limit theorem concerns weak convergence of distribution
measures:

\[
\boxed{
P_{
(S_n-n\mu)/(\sigma\sqrt{n})
}
\Rightarrow
N(0,1).
}
\]

Thus the LLN and CLT use different structures:

\[
\boxed{
\text{LLN}
\longrightarrow
\text{sample-path convergence},
}
\]

\[
\boxed{
\text{CLT}
\longrightarrow
\text{distributional convergence}.
}
\]

%%%%%%%%%%%%%%%%%%%%%%%%%%%%%%%%%%%%%%%%%%%%%%%%%%%%%%%%%%%%
\subsection{Why This Framework Matters for Stochastic Processes}
\label{subsec:measure-framework-stochastic-processes}
%%%%%%%%%%%%%%%%%%%%%%%%%%%%%%%%%%%%%%%%%%%%%%%%%%%%%%%%%%%%

A stochastic process may involve an uncountable collection

\[
\{X_t:t\in T\}
\]

of random variables.

Questions then involve:

\[
\text{joint measurability},
\]

\[
\text{filtrations},
\]

\[
\text{conditional expectations},
\]

\[
\text{stopping times},
\]

and

\[
\text{path properties}.
\]

These concepts cannot be handled adequately without measure theory.

%%%%%%%%%%%%%%%%%%%%%%%%%%%%%%%%%%%%%%%%%%%%%%%%%%%%%%%%%%%%
\subsection{Probability Measure Versus Density}
\label{subsec:measure-versus-density}
%%%%%%%%%%%%%%%%%%%%%%%%%%%%%%%%%%%%%%%%%%%%%%%%%%%%%%%%%%%%

A probability density is not the probability distribution itself.

The probability measure is the fundamental object.

If a density exists,

\[
\boxed{
P_X(B)
=
\int_B
f_X\,d\lambda.
}
\]

But some distributions have no density.

Thus:

\[
\boxed{
\text{probability measure}
\text{ is more general than }
\text{probability density}.
}
\]

%%%%%%%%%%%%%%%%%%%%%%%%%%%%%%%%%%%%%%%%%%%%%%%%%%%%%%%%%%%%
\subsection{Common Errors}
\label{subsec:measure-probability-common-errors}
%%%%%%%%%%%%%%%%%%%%%%%%%%%%%%%%%%%%%%%%%%%%%%%%%%%%%%%%%%%%

\subsubsection{Assuming Every Subset Is an Event}

On uncountable spaces, not every subset need be measurable.

Probabilities are defined only for events in the sigma-algebra.

%%%%%%%%%%%%%%%%%%%%%%%%%%%%%%%%%%%%%%%%%%%%%%%%%%%%%%%%%%%%
\subsubsection{Confusing the Sample Space with the Sigma-Algebra}

The sample space is

\[
\Omega.
\]

The sigma-algebra is a collection of subsets:

\[
\mathcal{F}\subseteq2^\Omega.
\]

They are different objects.

%%%%%%%%%%%%%%%%%%%%%%%%%%%%%%%%%%%%%%%%%%%%%%%%%%%%%%%%%%%%
\subsubsection{Confusing a Random Variable with Its Distribution}

A random variable is a measurable function.

Its distribution is a probability measure induced by that function.

Different random variables can have the same law.

%%%%%%%%%%%%%%%%%%%%%%%%%%%%%%%%%%%%%%%%%%%%%%%%%%%%%%%%%%%%
\subsubsection{Assuming Every Continuous Distribution Has a PDF}

A continuous CDF need not be absolutely continuous.

Singular continuous distributions exist.

%%%%%%%%%%%%%%%%%%%%%%%%%%%%%%%%%%%%%%%%%%%%%%%%%%%%%%%%%%%%
\subsubsection{Ignoring Almost-Sure Equality}

In \(L^p\) spaces, random variables that differ only on probability-zero
sets are treated as equivalent.

%%%%%%%%%%%%%%%%%%%%%%%%%%%%%%%%%%%%%%%%%%%%%%%%%%%%%%%%%%%%
\subsubsection{Interchanging Limits and Expectations Without a Theorem}

Almost-sure convergence alone does not imply

\[
\mathbb{E}[X_n]\to\mathbb{E}[X].
\]

A result such as monotone convergence, dominated convergence, or uniform
integrability is needed.

%%%%%%%%%%%%%%%%%%%%%%%%%%%%%%%%%%%%%%%%%%%%%%%%%%%%%%%%%%%%
\subsubsection{Confusing Conditional Expectation with a Number}

The quantity

\[
\mathbb{E}[X\mid\mathcal{G}]
\]

is generally a random variable, not a constant.

%%%%%%%%%%%%%%%%%%%%%%%%%%%%%%%%%%%%%%%%%%%%%%%%%%%%%%%%%%%%
\subsubsection{Treating Conditioning on \(Y=y\) as an Elementary Ratio}

For continuous \(Y\),

\[
P(Y=y)=0.
\]

Conditional distributions require densities, kernels, or regular
conditional probabilities rather than the naive ratio formula.

%%%%%%%%%%%%%%%%%%%%%%%%%%%%%%%%%%%%%%%%%%%%%%%%%%%%%%%%%%%%
\subsubsection{Confusing a.s. Convergence and \(L^p\) Convergence}

Neither implies the other in complete generality without additional
assumptions.

%%%%%%%%%%%%%%%%%%%%%%%%%%%%%%%%%%%%%%%%%%%%%%%%%%%%%%%%%%%%
\subsubsection{Assuming Independence Is Only an Event-Level Concept}

Modern probability defines independence for sigma-algebras, random
variables, random vectors, and stochastic processes.

%%%%%%%%%%%%%%%%%%%%%%%%%%%%%%%%%%%%%%%%%%%%%%%%%%%%%%%%%%%%
\subsection{Applications}
\label{subsec:measure-probability-applications}
%%%%%%%%%%%%%%%%%%%%%%%%%%%%%%%%%%%%%%%%%%%%%%%%%%%%%%%%%%%%

\begin{applicationbox}

\textbf{Stochastic Processes.}
Filtrations, conditional expectations, path spaces, and transition
kernels rely on measure theory.

\medskip

\textbf{Stochastic Calculus.}
Brownian motion, It\^o integrals, martingales, and stochastic
differential equations require measure-theoretic probability.

\medskip

\textbf{Statistics.}
Conditional expectation, likelihood, asymptotic convergence, Bayesian
models, and probability kernels depend on the measure framework.

\medskip

\textbf{Machine Learning.}
Probabilistic models, latent variables, posterior distributions, and
expectations over high-dimensional spaces use measurable-space
structures.

\medskip

\textbf{Finance.}
Risk-neutral measures, conditional expectations, martingales, and
derivative pricing rely directly on probability measures.

\medskip

\textbf{Ergodic Theory.}
Long-run averages are studied using measure-preserving transformations
and almost-sure convergence.

\medskip

\textbf{Functional Analysis.}
\(L^p\) spaces of random variables connect probability with Banach and
Hilbert space theory.

\medskip

\textbf{Modern Probability.}
Nearly every advanced theorem in probability is naturally formulated on
a measure space.

\end{applicationbox}

%%%%%%%%%%%%%%%%%%%%%%%%%%%%%%%%%%%%%%%%%%%%%%%%%%%%%%%%%%%%
\subsection{Exercises}
\label{subsec:measure-probability-exercises}
%%%%%%%%%%%%%%%%%%%%%%%%%%%%%%%%%%%%%%%%%%%%%%%%%%%%%%%%%%%%

\begin{exercise}[1]
Define a probability space.
\end{exercise}

\begin{exercise}[1]
Define a sigma-algebra.
\end{exercise}

\begin{exercise}[1]
Define a probability measure.
\end{exercise}

\begin{exercise}[1]
What is a measurable event?
\end{exercise}

\begin{exercise}[1]
Define a real-valued random variable measure-theoretically.
\end{exercise}

\begin{exercise}[1]
Define the distribution of a random variable as a pushforward measure.
\end{exercise}

\begin{exercise}[1]
Define an integrable random variable.
\end{exercise}

\begin{exercise}[1]
Define convergence in \(L^p\).
\end{exercise}

\begin{exercise}[1]
Define conditional expectation with respect to a sigma-algebra.
\end{exercise}

\begin{exercise}[1]
Define independent sigma-algebras.
\end{exercise}

\begin{exercise}[2]
Show that every sigma-algebra contains the empty set.
\end{exercise}

\begin{exercise}[2]
Show that sigma-algebras are closed under countable intersections.
\end{exercise}

\begin{exercise}[2]
Use countable additivity to prove

\[
P(\varnothing)=0.
\]
\end{exercise}

\begin{exercise}[2]
Prove monotonicity:

\[
A\subseteq B
\Longrightarrow
P(A)\leq P(B).
\]
\end{exercise}

\begin{exercise}[2]
Prove

\[
P(A^c)=1-P(A).
\]
\end{exercise}

\begin{exercise}[2]
Prove continuity from below.
\end{exercise}

\begin{exercise}[2]
Prove continuity from above.
\end{exercise}

\begin{exercise}[2]
Show that

\[
\mathbf{1}_A
\]

is measurable whenever \(A\in\mathcal{F}\).
\end{exercise}

\begin{exercise}[2]
Show that

\[
\mathbb{E}[\mathbf{1}_A]
=
P(A).
\]
\end{exercise}

\begin{exercise}[2]
Write a finite-valued random variable as a linear combination of
indicator functions.
\end{exercise}

\begin{exercise}[2]
Let

\[
X_n=\min(X,n)
\]

for \(X\geq0\). Show that

\[
X_n\uparrow X.
\]
\end{exercise}

\begin{exercise}[2]
Use monotone convergence to show

\[
\mathbb{E}[\min(X,n)]
\to
\mathbb{E}[X].
\]
\end{exercise}

\begin{exercise}[3]
Prove the tail-integral formula

\[
\mathbb{E}[X]
=
\int_0^\infty
P(X>t)\,dt
\]

for nonnegative \(X\) using Tonelli's theorem.
\end{exercise}

\begin{exercise}[3]
Show that if

\[
X_n\xrightarrow{L^p}X,
\]

then

\[
X_n\xrightarrow{P}X.
\]
\end{exercise}

\begin{exercise}[3]
Give an example of almost-sure convergence that does not imply
\(L^1\) convergence.
\end{exercise}

\begin{exercise}[3]
Give an example of convergence in probability that does not imply
almost-sure convergence along the full sequence.
\end{exercise}

\begin{exercise}[3]
Prove that if

\[
X_n\xrightarrow{\text{a.s.}}X
\]

and

\[
|X_n|\leq Y
\]

for an integrable \(Y\), then

\[
X_n\xrightarrow{L^1}X.
\]
\end{exercise}

\begin{exercise}[3]
Prove the tower property for a simple finite conditioning sigma-algebra.
\end{exercise}

\begin{exercise}[3]
Show that if \(X\) is \(\mathcal{G}\)-measurable, then

\[
\mathbb{E}[X\mid\mathcal{G}]
=
X
\]

almost surely.
\end{exercise}

\begin{exercise}[3]
Show that if \(X\) is independent of \(\mathcal{G}\), then

\[
\mathbb{E}[X\mid\mathcal{G}]
=
\mathbb{E}[X].
\]
\end{exercise}

\begin{exercise}[3]
Show that conditional probability can be written as

\[
P(A\mid\mathcal{G})
=
\mathbb{E}
[
\mathbf{1}_A
\mid
\mathcal{G}
].
\]
\end{exercise}

\begin{exercise}[3]
Show that independence of random variables is equivalent to independence
of the sigma-algebras they generate.
\end{exercise}

\begin{exercise}[3]
Use Hölder's inequality to derive Cauchy--Schwarz.
\end{exercise}

\begin{exercise}[3]
Use Cauchy--Schwarz to prove

\[
|\operatorname{Cov}(X,Y)|
\leq
\sigma_X\sigma_Y.
\]
\end{exercise}

\begin{exercise}[3]
Prove that conditional expectation is an orthogonal projection in
\(L^2\).
\end{exercise}

\begin{exercise}[3]
Use the projection interpretation to prove the law of total variance.
\end{exercise}

\begin{exercise}[3]
State and prove the first Borel--Cantelli lemma.
\end{exercise}

\begin{exercise}[3]
Explain why independence is not required for the first Borel--Cantelli
lemma.
\end{exercise}

\begin{exercise}[3]
State the second Borel--Cantelli lemma and identify where independence
is used.
\end{exercise}

\begin{exercise}[4]
Study the construction of Lebesgue integration from simple functions.
\end{exercise}

\begin{exercise}[4]
Study the Radon--Nikodym theorem and prove conditional expectation
existence using it.
\end{exercise}

\begin{exercise}[4]
Investigate singular continuous probability measures using the Cantor
distribution.
\end{exercise}

\begin{exercise}[4]
Study regular conditional probabilities on standard Borel spaces.
\end{exercise}

\begin{exercise}[4]
Investigate probability kernels and their composition.
\end{exercise}

\begin{exercise}[4]
Study the Kolmogorov extension theorem for constructing stochastic
processes.
\end{exercise}

\begin{exercise}[4]
Study the Portmanteau theorem for weak convergence of probability
measures.
\end{exercise}

\begin{exercise}[4]
Investigate Prokhorov's theorem and the role of tightness.
\end{exercise}

\begin{exercise}[4]
Study uniform integrability and prove a version of Vitali's convergence
theorem.
\end{exercise}

\begin{exercise}[4]
Investigate martingales as sequences defined by conditional
expectations.
\end{exercise}

\begin{exercise}[4]
Study filtrations, stopping times, and optional sampling.
\end{exercise}

\begin{exercise}[4]
Investigate \(L^p\) spaces as Banach spaces and \(L^2\) as a Hilbert
space.
\end{exercise}

\begin{exercise}[4]
Study infinite product probability spaces and independent coordinate
random variables.
\end{exercise}

%%%%%%%%%%%%%%%%%%%%%%%%%%%%%%%%%%%%%%%%%%%%%%%%%%%%%%%%%%%%
\subsection{Conceptual Summary}
\label{subsec:measure-probability-summary}
%%%%%%%%%%%%%%%%%%%%%%%%%%%%%%%%%%%%%%%%%%%%%%%%%%%%%%%%%%%%

A probability space is

\[
\boxed{
(\Omega,\mathcal{F},P).
}
\]

The sigma-algebra \(\mathcal{F}\) specifies measurable events.

The probability measure satisfies

\[
\boxed{
P(\Omega)=1
}
\]

and countable additivity.

A random variable is a measurable function

\[
\boxed{
X:
(\Omega,\mathcal{F})
\to
(\mathbb{R},\mathcal{B}(\mathbb{R})).
}
\]

Its law is the pushforward measure

\[
\boxed{
P_X
=
P\circ X^{-1}.
}
\]

Expectation is the Lebesgue integral

\[
\boxed{
\mathbb{E}[X]
=
\int_\Omega X\,dP.
}
\]

For an event \(A\),

\[
\boxed{
P(A)
=
\mathbb{E}[\mathbf{1}_A].
}
\]

The major convergence theorems include

\[
\boxed{
\text{Monotone Convergence},
}
\]

\[
\boxed{
\text{Fatou's Lemma},
}
\]

and

\[
\boxed{
\text{Dominated Convergence}.
}
\]

Important convergence modes are

\[
\boxed{
\xrightarrow{\text{a.s.}},
\qquad
\xrightarrow{P},
\qquad
\xrightarrow{d},
\qquad
\xrightarrow{L^p}.
}
\]

Conditional expectation is characterized by

\[
\boxed{
\mathbb{E}[X\mid\mathcal{G}]
}
\]

being \(\mathcal{G}\)-measurable and satisfying

\[
\boxed{
\int_G
\mathbb{E}[X\mid\mathcal{G}]
\,dP
=
\int_G
X\,dP
}
\]

for every \(G\in\mathcal{G}\).

The tower property is

\[
\boxed{
\mathbb{E}
[
\mathbb{E}[X\mid\mathcal{G}]
]
=
\mathbb{E}[X].
}
\]

Conditional probability is

\[
\boxed{
P(A\mid\mathcal{G})
=
\mathbb{E}
[
\mathbf{1}_A
\mid
\mathcal{G}
].
}
\]

Measure-theoretic probability therefore provides the rigorous language
for

\[
\boxed{
\text{events}
+
\text{random variables}
+
\text{expectation}
+
\text{conditioning}
+
\text{convergence}
+
\text{stochastic processes}.
}
\]

%%%%%%%%%%%%%%%%%%%%%%%%%%%%%%%%%%%%%%%%%%%%%%%%%%%%%%%%%%%%
\subsection{Section Connections}
\label{subsec:measure-probability-connections}
%%%%%%%%%%%%%%%%%%%%%%%%%%%%%%%%%%%%%%%%%%%%%%%%%%%%%%%%%%%%

\chapterconnection{
Measure-Theoretic Probability provides the rigorous foundation beneath
all earlier sections of Chapter 7. Random variables become measurable
functions, distributions become pushforward measures, and expectations
become Lebesgue integrals. The convergence notions used in the Laws of
Large Numbers and Central Limit Theorems are formalized here, while
conditional expectation is extended from elementary formulas to
conditioning on sigma-algebras. The next section, Stochastic Processes,
uses filtrations, product spaces, conditional expectations, and
finite-dimensional distributions to study random systems evolving in
time. Stochastic Calculus and Stochastic Differential Equations later
depend heavily on the martingale and filtration framework introduced
here.
}

%%%%%%%%%%%%%%%%%%%%%%%%%%%%%%%%%%%%%%%%%%%%%%%%%%%%%%%%%%%%
% SECTION SOURCE TRACKING
%%%%%%%%%%%%%%%%%%%%%%%%%%%%%%%%%%%%%%%%%%%%%%%%%%%%%%%%%%%%

%------------------------------------------------------------
% 7.11 Measure-Theoretic Probability
%------------------------------------------------------------
%
% Source basis:
%   - Standard measure theory.
%   - Standard measure-theoretic probability.
%   - Standard Lebesgue integration and conditional expectation theory.
%
% Used for:
%   - probability spaces;
%   - sigma-algebras;
%   - measurable events;
%   - generated sigma-algebras;
%   - Borel sigma-algebra;
%   - probability measures;
%   - Kolmogorov axioms;
%   - continuity of probability measures;
%   - random variables as measurable functions;
%   - distributions as pushforward measures;
%   - discrete probability laws;
%   - absolutely continuous laws;
%   - Radon--Nikodym densities;
%   - singular distributions;
%   - indicator functions;
%   - simple random variables;
%   - expectation as a Lebesgue integral;
%   - positive and negative parts;
%   - integrability;
%   - almost-sure language;
%   - monotone convergence theorem;
%   - Fatou's lemma;
%   - dominated convergence theorem;
%   - Tonelli and Fubini methods;
%   - tail-integral formulas;
%   - Lp spaces;
%   - Lp convergence;
%   - convergence modes;
%   - sigma-algebras as information;
%   - conditional expectation;
%   - conditional probability;
%   - tower property;
%   - projection interpretation;
%   - law of total variance;
%   - Radon--Nikodym theorem;
%   - independence of sigma-algebras;
%   - product probability spaces;
%   - Borel--Cantelli lemmas;
%   - regular conditional probability preview;
%   - probability kernels preview;
%   - filtrations preview;
%   - martingale preview;
%   - complete probability spaces;
%   - Jensen, Hölder, Cauchy--Schwarz, and Minkowski inequalities;
%   - uniform integrability preview;
%   - weak convergence and tightness preview;
%   - applications and exercises.
%
% Notes:
%   - Stochastic Processes are developed next in Section 7.12.
%   - Martingale theory is expanded within stochastic-process and
%     stochastic-calculus material.
%   - General measure theory and Lebesgue integration were developed in
%     Chapter 5 and are used here in probabilistic form.
%
%%%%%%%%%%%%%%%%%%%%%%%%%%%%%%%%%%%%%%%%%%%%%%%%%%%%%%%%%%%%